
This study validates that published technical reports from three major LLM labs (OpenAI, Anthropic, Meta) contain category-level error rate data for TruthfulQA and MMLU across model generations. Data availability meets predefined criteria: 3/3 family coverage, 12-15 categories per benchmark exceeding the 10-category threshold, 100\% completeness using manual extraction, and temporal coverage across both baseline and current models.

These findings establish that category-level benchmark data are systematically available in published materials. Whether such data support specific applications (error taxonomy generation, pattern analysis, weak supervision approaches) remains to be determined. Validating those applications would require:

\begin{enumerate}
\item Testing whether metadata features correlate with category-level error rates
\item Testing whether category-level data enable item-level clustering with adequate variance separation
\item Testing whether resulting clusters are interpretable via expert agreement (e.g., Cohen's kappa ≥0.7)
\item Testing whether patterns generalize across benchmarks
\end{enumerate}

Each of these represents a distinct research question beyond data availability validation. The current study establishes only that the data prerequisite is satisfied; the mechanism and applications remain unvalidated.

\textbf{Future work.} Three directions could build on these findings: (1) developing automated extraction systems to enable scalable tracking across more families and longer time periods, achieving target success rates of 85-95\%, (2) validating specific downstream applications such as taxonomy generation with formal evaluation of effectiveness, and (3) investigating temporal stability of error patterns to determine whether categories reflect enduring limitations or transient artifacts.
